\documentclass[11pt,a4paper]{report}

% ---- Packages (all common; compiles with pdflatex) ----------------------
\usepackage[utf8]{inputenc}
\usepackage[T1]{fontenc}
\usepackage[margin=1in]{geometry}
\usepackage{parskip}
\usepackage{xcolor}
\usepackage{listings}
\usepackage{enumitem}
\usepackage{longtable}
\usepackage{titlesec}
\usepackage[hidelinks]{hyperref}

% ---- Listings style ------------------------------------------------------
\definecolor{codebg}{RGB}{245,246,249}
\definecolor{codekw}{RGB}{37,99,235}
\definecolor{codecomment}{RGB}{107,114,128}
\definecolor{codestr}{RGB}{21,128,61}

\lstset{
  basicstyle=\ttfamily\footnotesize,
  backgroundcolor=\color{codebg},
  keywordstyle=\color{codekw}\bfseries,
  commentstyle=\color{codecomment}\itshape,
  stringstyle=\color{codestr},
  showstringspaces=false,
  breaklines=true,
  columns=fullflexible,
  frame=single,
  framesep=4pt,
  rulecolor=\color{gray!30},
  numbers=none,
  tabsize=2,
  aboveskip=8pt,
  belowskip=8pt,
}

% Mock-interview question macro. Escape underscores as \_ inside arguments;
% avoid raw &, %, #, $ (use words).
\newcommand{\mq}[5]{%
  \par\smallskip\noindent\textbf{Q. #1}\par
  \textit{Ideal:} #2\par
  \textit{Hints:} #3\par
  \textit{Common mistakes:} #4\par
  \textit{Follow-up:} #5\par\smallskip
}

\titleformat{\chapter}[display]{\normalfont\huge\bfseries}{\chaptertitlename\ \thechapter}{20pt}{\Huge}

\title{\textbf{RBAC Service --- Feature Extensions \\ and Interview Preparation Guide}\\[6pt]
\large Companion to the original submission report (\texttt{report.tex})}
\author{sec-task-2 (interview playground)}
\date{}

\begin{document}

\maketitle

\begin{center}
\textit{This guide builds on the original \texttt{report.tex}. It does not repeat
the basics of RBAC, FastAPI, or the core permission check --- see that report
for those. Here the focus is: the follow-up features actually implemented in
\texttt{sec-task-2}, a repeatable way to design any new feature on the spot, the
SQL migration, and a large bank of practice questions and exercises.}
\end{center}

\tableofcontents

% =========================================================================
\chapter{Why Interviewers Ask Follow-up Implementation Questions}
% =========================================================================

After a take-home, interviewers rarely re-check whether your original code works.
Instead they say ``now add one more feature'' or ``how would you change this to
support X''. The point is to watch you \emph{work}, live, on code you understand.

\section{What They Are Really Evaluating}

\begin{itemize}[leftmargin=1.4em]
  \item \textbf{Coding fluency.} Can you turn an idea into correct code without
        fighting the language or the framework?
  \item \textbf{Architecture instinct.} Do your changes respect the existing
        layering (routes \texttt{->} service \texttt{->} storage), or do you hack
        logic into a route handler?
  \item \textbf{Thinking process.} Do you clarify the requirement, consider edge
        cases, and reason about complexity before typing?
  \item \textbf{Debugging.} When something breaks, do you read the error, form a
        hypothesis, and test it --- or guess randomly?
  \item \textbf{API design.} Do you pick the right verb, path, status code, and
        response shape?
  \item \textbf{Communication.} Do you narrate what you are doing and why, and ask
        good questions?
\end{itemize}

\section{How to Score Well}

Talk before you type. State your plan in one or two sentences (``I'll add a
service method, a route, a schema, and two tests''). Name the edge cases out
loud. Keep each change small and runnable. Re-run the tests often. If you get
stuck, say what you expected versus what happened. Interviewers reward a calm,
structured approach far more than raw speed.

\section{This Repo as Evidence}

Everything discussed here is real and runnable in \texttt{sec-task-2}: an
extended in-memory backend (\texttt{app/}), a parallel stdlib \texttt{sqlite3}
version (\texttt{sql\_version/}), and a full test suite. When an interviewer asks
``implement delete role with cascade'', you can point to exactly how it was done
in both storage backends.

% =========================================================================
\chapter{Implemented Features: Deep Dive}
% =========================================================================

For each feature: why it is useful, which files change and why, the thinking, the
actual code, complexity, edge cases, common bugs, likely follow-ups, and
alternatives. All layer names refer to routes (\texttt{routes.py}), service
(\texttt{services.py}), storage (\texttt{storage.py}), and schemas
(\texttt{schemas.py}).

\section{Delete User / Delete Role (with cascade)}
\textbf{Why useful.} Real systems remove people and roles; leaving them forever
is a security and hygiene problem.\\
\textbf{Files.} \texttt{services.py} (logic), \texttt{routes.py} (DELETE
endpoints). \textbf{Why:} deletion is a state change (service) exposed over HTTP
(route).\\
\textbf{Thinking.} Deleting a user is a simple pop. Deleting a role is harder:
users reference role ids, so a naive delete leaves \emph{dangling references}.
The role delete must cascade --- remove the id from every user.\\
\textbf{Actual code (in-memory).}
\begin{lstlisting}[language=Python]
def delete_user(self, user_id):
    self.get_user(user_id)          # 404 if missing
    del self.store.users[user_id]
    self._invalidate_cache(user_id)

def delete_role(self, role_id):
    self.get_role(role_id)          # 404 if missing
    del self.store.roles[role_id]
    for user in self.store.users.values():
        user.role_ids.discard(role_id)   # cascade
    self._invalidate_cache()
\end{lstlisting}
\textbf{SQL version.} \texttt{conn.execute("DELETE FROM roles WHERE id=?", (rid,))}
--- \texttt{ON DELETE CASCADE} removes the join-table rows for that role
automatically (each connection sets \texttt{PRAGMA foreign\_keys = ON}), so no
dangling references.\\
\textbf{Complexity.} \texttt{delete\_user} O(1); \texttt{delete\_role} O(U) in
the number of users (the cascade loop).\\
\textbf{Edge cases.} Deleting a non-existent id (404); deleting a role currently
in use (must cascade); double delete (second one 404).\\
\textbf{Common bugs.} Forgetting the cascade (dangling ids); forgetting to
invalidate the permission cache.\\
\textbf{Follow-ups.} Soft delete instead of hard delete? Return 204 or the
deleted object? What about audit?\\
\textbf{Alternatives.} Soft delete (\texttt{deleted\_at}); block deletion of a
role still assigned (409) instead of cascading.

\section{Get User / Role by Id}
\textbf{Why useful.} Clients need to fetch a single resource, not just lists.\\
\textbf{Files.} \texttt{routes.py} (the service already had \texttt{get\_user} /
\texttt{get\_role}).\\
\textbf{Code.}
\begin{lstlisting}[language=Python]
@router.get("/users/{user_id}", response_model=UserResponse)
def get_user(user_id: str, service = Depends(get_service)):
    return _user_to_response(service.get_user(user_id))
\end{lstlisting}
\textbf{Complexity.} O(1) dict lookup. \textbf{Edge case:} unknown id -> 404.\\
\textbf{Common bug.} Route ordering: \texttt{/users/count} must be declared before
\texttt{/users/\{user\_id\}} or ``count'' is captured as an id.

\section{Rename User / Rename Role (PATCH)}
\textbf{Why useful.} Names change; you should not have to delete and recreate.\\
\textbf{Files.} \texttt{schemas.py} (\texttt{RenameRequest}), \texttt{services.py}
(\texttt{rename\_user} / \texttt{rename\_role}), \texttt{routes.py} (PATCH).\\
\textbf{Thinking.} PATCH = partial update. Must trim, reject blanks, and enforce
uniqueness \emph{excluding the current record}, and bump \texttt{updated\_at}.\\
\textbf{Code.}
\begin{lstlisting}[language=Python]
def rename_user(self, user_id, new_name):
    user = self.get_user(user_id)
    cleaned = clean_name(new_name)
    self._assert_unique_user_name(cleaned, exclude_id=user_id)
    user.name = cleaned
    user.updated_at = utcnow()
    return user
\end{lstlisting}
\textbf{Complexity.} O(N) for the uniqueness scan in memory (O(1) in SQL with a
unique index). \textbf{Edge cases:} rename to same name (allowed, exclude self);
rename to another's name (409); blank name (422).\\
\textbf{Common bug.} Uniqueness check that does not exclude the record being
renamed, so renaming to its own name wrongly conflicts.\\
\textbf{Follow-up.} Why PATCH not PUT? (PATCH = partial; PUT = full replace.)

\section{Replace User Roles / Replace Role Permissions (PUT)}
\textbf{Why useful.} Setting the full list at once is common (an admin UI toggles
checkboxes and saves).\\
\textbf{Files.} \texttt{schemas.py}, \texttt{services.py}, \texttt{routes.py}.\\
\textbf{Thinking.} PUT = replace the whole collection. Validate \emph{all}
referenced ids first (all-or-nothing) before mutating, so a bad id does not leave
a half-applied state.\\
\textbf{Code.}
\begin{lstlisting}[language=Python]
def replace_user_roles(self, user_id, role_ids):
    user = self.get_user(user_id)
    for role_id in role_ids:
        self.get_role(role_id)      # validate ALL before assigning
    user.role_ids = set(role_ids)
    user.updated_at = utcnow()
    self._invalidate_cache(user_id)
    return user
\end{lstlisting}
\textbf{Complexity.} O(K) in the number of ids supplied. \textbf{Edge cases:}
empty list (clears roles); duplicate ids (set de-dups); unknown id (404, nothing
changed).\\
\textbf{Common bug.} Assigning as you validate, so a later bad id leaves a
partial update.

\section{List a User's Effective Permissions}
\textbf{Why useful.} ``What can this user actually do?'' is a core RBAC question
and powers UIs and debugging.\\
\textbf{Files.} \texttt{services.py} (\texttt{get\_effective\_permissions}),
\texttt{routes.py}.\\
\textbf{Thinking.} A user's effective permissions are the \emph{union} of all
their roles' permission sets. This is also the perfect thing to cache.\\
\textbf{Code.}
\begin{lstlisting}[language=Python]
def get_effective_permissions(self, user_id):
    if user_id in self.store.permission_cache:
        return self.store.permission_cache[user_id]
    user = self.get_user(user_id)
    permissions = set()
    for role_id in user.role_ids:
        role = self.store.roles.get(role_id)
        if role is not None:
            permissions |= role.permissions
    self.store.permission_cache[user_id] = permissions
    return permissions
\end{lstlisting}
\textbf{SQL version.} A single \texttt{DISTINCT} JOIN across
\texttt{user\_roles} and \texttt{role\_permissions}.\\
\textbf{Complexity.} O(R*P) to compute (roles times permissions), O(1) on a cache
hit. \textbf{Edge cases:} user with no roles (empty set); unknown user (404).

\section{List Users Having a Role}
\textbf{Why useful.} ``Who is an admin?'' --- audits and impact analysis before
deleting a role.\\
\textbf{Files.} \texttt{services.py}, \texttt{routes.py}
(\texttt{GET /roles/\{role\_id\}/users}).\\
\textbf{Code.}
\begin{lstlisting}[language=Python]
def list_users_with_role(self, role_id):
    self.get_role(role_id)          # 404 if the role is missing
    return [u for u in self.store.users.values() if role_id in u.role_ids]
\end{lstlisting}
\textbf{Complexity.} O(U) scan in memory; in SQL a JOIN on \texttt{user\_roles}
(indexed).\\
\textbf{Follow-up.} At scale, keep a reverse index (role -> user ids) to avoid the
scan.

\section{Search Users / Roles by Name}
\textbf{Why useful.} Lists get long; users expect to filter.\\
\textbf{Files.} \texttt{services.py} (\texttt{search\_users}), \texttt{routes.py}
(query param \texttt{name}).\\
\textbf{Thinking.} A query parameter, not a new path --- it filters the same
collection. Case-insensitive substring match.\\
\textbf{Code.}
\begin{lstlisting}[language=Python]
@router.get("/users", response_model=list[UserResponse])
def list_users(name: str | None = Query(default=None), service = Depends(get_service)):
    return [_user_to_response(u) for u in service.search_users(name)]
\end{lstlisting}
\textbf{Complexity.} O(N) scan in memory; in SQL a \texttt{LIKE} (add a trigram or
prefix index for large tables). \textbf{Edge cases:} empty query returns all;
whitespace query trimmed.\\
\textbf{Follow-up.} Add pagination (\texttt{limit}/\texttt{offset}); sort order.

\section{Duplicate Prevention and Validation / Normalization}
\textbf{Why useful.} Prevents confusing data (two ``admin'' roles) and garbage
(blank names, mixed-case permissions that never match).\\
\textbf{Files.} \texttt{services.py} (\texttt{clean\_name},
\texttt{clean\_permission}, uniqueness asserts), \texttt{errors.py}
(\texttt{ValidationError}, \texttt{ConflictError}).\\
\textbf{Code.}
\begin{lstlisting}[language=Python]
def clean_name(name):
    cleaned = (name or "").strip()
    if not cleaned:
        raise ValidationError("Name must not be blank")
    return cleaned

def clean_permission(resource, action):
    r = (resource or "").strip().lower()
    a = (action or "").strip().lower()
    if not r or not a:
        raise ValidationError("Permission resource and action must not be blank")
    return Permission(resource=r, action=a)
\end{lstlisting}
\textbf{Thinking.} Normalize at the boundary (trim names, lowercase
permissions) so comparisons are predictable; enforce case-insensitive uniqueness
excluding self on rename.\\
\textbf{Edge cases.} Whitespace-only name (\texttt{"   "}) passes Pydantic
\texttt{min\_length=1} but is caught by \texttt{clean\_name} -> 422; ``Admin''
versus ``admin'' collide (409).\\
\textbf{Common bug.} Storing permissions with mixed case so
\texttt{documents:Write} never matches \texttt{documents:write}.

\section{Centralized Error Responses}
\textbf{Why useful.} Consistent, machine-readable errors are easier for clients
and cleaner than try/except in every route.\\
\textbf{Files.} \texttt{errors.py} (exception hierarchy), \texttt{main.py}
(handlers).\\
\textbf{Code.}
\begin{lstlisting}[language=Python]
class AppError(Exception):
    status_code = 400
    code = "error"
    def __init__(self, message):
        self.message = message
        super().__init__(message)

@app.exception_handler(AppError)
def handle_app_error(request, exc):
    return JSONResponse(status_code=exc.status_code,
                        content={"error": {"code": exc.code, "message": exc.message}})
\end{lstlisting}
\textbf{Result.} Every failure looks like
\texttt{\{"error": \{"code": "not\_found", "message": "..."\}\}}. A second handler
maps Pydantic validation errors to the same shape with code
\texttt{validation\_error}.\\
\textbf{Follow-up.} Add a request id / correlation id to every error for tracing.

\section{Effective-Permission Cache (with invalidation)}
\textbf{Why useful.} Permission checks are the read-hot path; caching the union
per user makes a check an O(1) lookup.\\
\textbf{Files.} \texttt{storage.py} (\texttt{permission\_cache}),
\texttt{services.py} (read-through + \texttt{\_invalidate\_cache} on writes).\\
\textbf{Thinking.} The hard part is invalidation: user role change -> invalidate
that user; role permission change or role delete -> invalidate broadly.\\
\textbf{Common bug.} Forgetting to invalidate, so a revoked permission still
appears granted (this repo has a test that fails without invalidation).\\
\textbf{Follow-up.} Move the cache to Redis for multi-process; add a TTL as a
safety net.

\section{Timestamps, Audit Log, Health, Version, Stats (bonus)}
\textbf{Why useful.} Operability: know when things changed, what happened, whether
the service is up, and how big the data is.\\
\textbf{Files.} \texttt{models.py} (\texttt{created\_at} / \texttt{updated\_at}),
\texttt{storage.py} (\texttt{audit\_log}), \texttt{services.py}
(\texttt{\_log}, counts, \texttt{get\_audit\_log}), \texttt{routes.py}
(\texttt{/health}, \texttt{/system/info}, \texttt{/system/audit},
\texttt{/users/count}, \texttt{/roles/count}).\\
\textbf{Edge cases.} \texttt{updated\_at} should change only on real
modifications; \texttt{/health} must not touch storage so it stays fast and
always answers.\\
\textbf{Follow-up.} Persist audit to a durable store; add readiness vs liveness
probes; expose Prometheus metrics.

% =========================================================================
\chapter{Feature Design Framework}
% =========================================================================

When an interviewer says ``implement X'', run this pipeline out loud. It works for
almost any backend feature.

\begin{verbatim}
  Requirement  -> clarify scope, inputs, outputs, who calls it
       |
  Affected     -> which layers? (schema / route / service / storage)
   layer
       |
  Storage      -> new data? new field? new index? migration?
   impact
       |
  Service      -> the rule/algorithm; validation; normalization
   impact
       |
  Route        -> verb, path, status code, request + response schema
   impact
       |
  Schema       -> request/response Pydantic models
   impact
       |
  Testing      -> success, failure, edge, duplicate, validation
       |
  Complexity   -> time/space; will it scale?
       |
  Edge cases   -> empty, missing, duplicate, concurrent, large
       |
  Validation   -> trim, required, ranges, referential integrity
       |
  Response     -> shape, status, errors, idempotency
\end{verbatim}

\section{Reusable Checklists}

\textbf{New write endpoint checklist}
\begin{enumerate}[leftmargin=1.5em]
  \item Pick the verb: create=POST, full replace=PUT, partial=PATCH,
        remove=DELETE.
  \item Define request schema (validated) and response schema.
  \item Add a service method that validates, normalizes, mutates, and returns.
  \item Validate references (do the ids exist? 404) and uniqueness (409).
  \item Pick the status code (201 create, 200 update, 204 delete).
  \item Invalidate caches / update timestamps / write audit if relevant.
  \item Tests: happy path, not-found, validation error, duplicate, idempotency.
\end{enumerate}

\textbf{New read endpoint checklist}
\begin{enumerate}[leftmargin=1.5em]
  \item Single resource (path param) or collection (maybe filters/pagination)?
  \item Declare static paths before parametrized ones.
  \item 404 for a missing single resource; empty list is fine for a collection.
  \item Response schema; sort for deterministic output.
  \item Tests: found, not-found, filter matches/misses.
\end{enumerate}

\textbf{``Does this scale?'' checklist}
\begin{enumerate}[leftmargin=1.5em]
  \item What is the complexity? Any O(n) scan on the hot path?
  \item Is it the read path (cache it) or the write path (keep correct)?
  \item Does it need an index? A reverse index? Pagination?
  \item Is state per-process (breaks horizontal scaling) or shared?
\end{enumerate}

% =========================================================================
\chapter{On-the-Spot Coding Strategy}
% =========================================================================

\section{Think Before You Type}
Restate the task in your own words. Confirm inputs, outputs, and the one or two
edge cases that matter. Sketch the 3--4 changes (schema, service, route, test)
before writing any of them.

\section{Communicate}
Narrate: ``I'll put the rule in the service so it's testable, add a thin route,
and cover not-found and duplicate cases.'' Silence reads as being stuck; a
running commentary reads as competence.

\section{Ask Clarifying Questions}
Good ones: ``Should delete be soft or hard?'', ``Is renaming to an existing name
allowed?'', ``Do you want pagination now or later?'' They show product sense and
prevent wasted work.

\section{Structure the Code}
Follow the existing layering. Small functions. Clear names. Put logic in the
service, keep routes thin. Reuse helpers (\texttt{clean\_name}, converters).

\section{Avoid Bugs}
Validate before you mutate (all-or-nothing). Handle the empty and missing cases
first. Re-run tests after each small step. Watch route ordering and cache
invalidation --- the two classic traps in this project.

\section{Finish Quickly}
Get a minimal correct version working, then improve. A working happy path plus one
edge case beats a half-written ``perfect'' design.

\section{Recover If Stuck}
Say what you expected versus what happened. Re-read the actual error. Shrink the
problem: hardcode a value, print state, or test the service method directly
without HTTP. Ask for a hint --- using one well is a positive signal.

% =========================================================================
\chapter{Backend Feature Catalog (approx. 150)}
% =========================================================================

Frequency = how often it comes up in interviews (H/M/L). Files use short names
(route = \texttt{routes.py}, svc = \texttt{services.py}, store =
\texttt{storage.py}, sch = \texttt{schemas.py}).

\section{Easy}
\begin{longtable}{p{0.30\textwidth} p{0.06\textwidth} p{0.56\textwidth}}
\textbf{Feature} & \textbf{Freq} & \textbf{Strategy (files)} \\
\hline
\endhead
Delete user & H & pop from store; route DELETE (svc, route) \\
Delete role with cascade & H & delete then discard id from all users (svc, route) \\
Get user by id & H & dict get, 404 if missing (route) \\
Get role by id & H & dict get, 404 (route) \\
List users & H & return all (route) \\
List roles & H & return all (route) \\
Rename user & M & trim, unique, update timestamp (svc, sch, route) \\
Rename role & M & same as rename user (svc, sch, route) \\
Count users & M & len of store (svc, route) \\
Count roles & M & len of store (svc, route) \\
Health endpoint & H & return status ok, no storage (route) \\
Version endpoint & M & return a constant version (route) \\
System info & M & name, version, counts (svc, route) \\
Assign role to user & H & add id to set (svc, route) \\
Remove role from user & H & discard id (svc, route) \\
Add permission to role & H & add to set (svc, route) \\
Remove permission from role & H & discard from set (svc, route) \\
Trim names & M & strip in clean helper (svc) \\
Reject blank names & M & raise validation (svc) \\
Lowercase permissions & M & normalize in clean helper (svc) \\
Duplicate user name check & H & scan/unique index, 409 (svc) \\
Duplicate role name check & H & same, 409 (svc) \\
Search users by name & H & query param, substring (svc, route) \\
Search roles by name & M & query param, substring (svc, route) \\
Users having a role & H & filter/JOIN, 404 if role missing (svc, route) \\
Effective permissions of a user & H & union of role perms (svc, route) \\
Replace user roles & M & validate all then set (svc, sch, route) \\
Replace role permissions & M & validate then set (svc, sch, route) \\
Return created object with 201 & M & set status\_code (route) \\
Return 204 on delete & M & no body (route) \\
Consistent error JSON & H & exception handler (errors, main) \\
Map validation to 422 & M & handler for RequestValidationError (main) \\
CORS for the frontend & M & add CORSMiddleware (main) \\
Serve static frontend & L & mount StaticFiles (main) \\
No-cache headers & L & middleware sets Cache-Control (main) \\
created\_at / updated\_at & M & timestamp fields, set on write (models, svc) \\
Idempotent assign (set) & M & set add is a no-op on repeat (svc) \\
Sort list output & L & sort before returning (route) \\
Reject empty permission & M & validate resource and action (svc) \\
Case-insensitive name uniqueness & M & compare lowercased (svc) \\
Get-or-create permission (SQL) & M & select then insert (svc) \\
Return effective perms sorted & L & sort tuples (route) \\
Simple audit log & M & append entries on writes (store, svc) \\
List audit entries & M & return tail of log (svc, route) \\
Whoami-style user detail & L & get by id plus roles (route) \\
Role detail with permissions & L & get by id plus perms (route) \\
Reject unknown role on create & H & validate ids, 404 (svc) \\
Clear a user's roles & L & replace with empty list (svc) \\
Ping/liveness & M & static ok (route) \\
Uptime in system info & L & now minus start time (route) \\
Echo request id header & L & read/return header (middleware) \\
Limit audit size & L & cap the list length (store) \\
Normalize whitespace in search & L & strip query (svc) \\
Return 409 body with code & M & conflict error shape (errors) \\
\end{longtable}

\section{Medium}
\begin{longtable}{p{0.30\textwidth} p{0.06\textwidth} p{0.56\textwidth}}
\textbf{Feature} & \textbf{Freq} & \textbf{Strategy (files)} \\
\hline
\endhead
Pagination (limit/offset) & H & slice list / SQL LIMIT OFFSET (svc, route) \\
Cursor pagination & M & keyset on id/created\_at (svc, route) \\
Sorting by field & M & sort key from query param (svc, route) \\
Multi-field filtering & M & combine query params (svc, route) \\
Bulk assign roles & M & loop validate then apply (svc, route) \\
Bulk create users & M & validate all, transaction (svc, route) \\
Soft delete & H & deleted\_at flag, filter reads (models, svc) \\
Restore soft-deleted & M & clear the flag (svc, route) \\
Wildcard permissions & H & match resource:* on check (svc) \\
Permission groups & M & named set expanded on assign (svc, models) \\
Effective-permission cache & H & memoize union, invalidate (store, svc) \\
Redis cache & M & external cache, invalidate on write (svc) \\
Optimistic locking & M & version column, 409 on stale (models, svc) \\
Transactions for multi-step & H & one unit of work commit (svc) \\
Request validation messages & M & Pydantic field errors (sch, main) \\
Rate limiting & M & token bucket per client (middleware) \\
Structured logging & M & logging module + middleware (main) \\
Request/correlation id & M & generate, thread through logs (middleware) \\
Export users/roles to JSON & M & serialize all (svc, route) \\
Import from JSON & M & validate then insert (svc, route) \\
Export to CSV & L & stream rows (route) \\
API key auth (simple) & M & header check dependency (deps) \\
Per-endpoint authorization & H & require-permission dependency (deps) \\
Role assignment expiry & M & store until-date, filter (models, svc) \\
Effective perms with groups & M & expand groups in union (svc) \\
Case-insensitive search index & M & functional/lower index (models) \\
Partial update of role & M & PATCH selected fields (sch, svc, route) \\
Reverse index role-to-users & M & maintain a map on writes (store) \\
Count users per role & M & group by / dict count (svc, route) \\
Bulk delete & M & validate all, one transaction (svc, route) \\
Idempotency keys & M & store key -> result (store, route) \\
ETag / If-Match & M & version-based concurrency (route) \\
Webhook on change & M & post event to a URL (svc) \\
Feature flag per tenant & M & flag map, check on request (svc) \\
Search with pagination & M & filter then slice (svc, route) \\
Deny list (explicit deny) & H & deny overrides allow in check (svc) \\
Time-boxed roles & M & valid-from/valid-to (models, svc) \\
Config via environment & M & read env, defaults (config) \\
Graceful shutdown & L & close resources on stop (main) \\
Batch permission check & M & check many at once (svc, route) \\
Cache stampede guard & M & lock/single-flight on miss (svc) \\
Audit filter by user/action & M & query params over log (svc, route) \\
Metrics endpoint & M & counters, latencies (route) \\
Retry-safe create & M & idempotency or unique key (svc) \\
Localization of errors & L & message catalog (errors) \\
Field-level permissions & M & permission per field (svc) \\
Pagination metadata & M & total, next cursor (route) \\
Read-only mode & L & reject writes when flag on (middleware) \\
Seed via CLI command & L & script calls services (scripts) \\
Validate permission vocabulary & M & allowed resources/actions (svc) \\
Dry-run mode for deletes & L & report impact, do not apply (route) \\
Bulk import with report & M & per-row success/failure (svc, route) \\
Rename with history & M & append old name to audit (svc) \\
\end{longtable}

\section{Hard}
\begin{longtable}{p{0.30\textwidth} p{0.06\textwidth} p{0.56\textwidth}}
\textbf{Feature} & \textbf{Freq} & \textbf{Strategy (files)} \\
\hline
\endhead
Role hierarchy (inheritance) & H & transitive closure, cycle guard (models, svc) \\
Permission inheritance cache & M & precompute closure, invalidate (svc) \\
Multi-tenant RBAC & H & tenant id scopes all queries (models, svc) \\
Attribute-based rules (ABAC) & M & policy over attributes (svc) \\
Delegated administration & M & who-can-grant-what rules (svc) \\
Just-in-time access & M & temporary grants with expiry (svc) \\
Separation of duties & M & conflicting-role constraints (svc) \\
Full audit with actor & H & who did what, immutable store (models, svc) \\
Event-driven invalidation & M & publish change events (svc, bus) \\
Read replicas routing & M & reads to replica, writes primary (db) \\
Sharding by tenant & L & route by tenant key (db) \\
Outbox pattern & M & write event in same transaction (svc, store) \\
Eventual-consistency cache & M & TTL plus event invalidation (svc) \\
Fine-grained resource ACLs & M & per-object permissions (models, svc) \\
Policy versioning & M & version rules, evaluate as-of (models) \\
Bulk recompute permissions & M & background job over users (worker) \\
Hierarchical resources & M & path-prefix permission match (svc) \\
Permission conditions & M & context predicate on grant (svc) \\
Concurrent write conflict resolution & M & optimistic version + retry (svc) \\
Migration with zero downtime & M & dual-write/read then cutover (db) \\
GDPR delete (erase user data) & M & cascade + anonymize audit (svc) \\
Approval workflow for grants & M & pending state + approve step (models, svc) \\
Cross-service authorization & M & central authz service/API (service) \\
Token-scoped permissions & M & scopes in JWT, intersect (deps) \\
Break-glass emergency access & L & audited elevated role (svc) \\
Permission simulation ("what if") & M & evaluate hypothetical roles (svc) \\
Blast-radius analysis on delete & M & compute affected users (svc) \\
Time-travel queries & L & history tables, as-of reads (db) \\
Rate limit per permission & L & bucket keyed by action (middleware) \\
Consistent hashing for cache & L & distribute keys across nodes (infra) \\
Materialized effective-perms view & M & maintained table, refresh (db) \\
Batch invalidation on role edit & H & find holders, invalidate all (svc) \\
Hierarchical role deletion & M & reparent or block children (svc) \\
Audit tamper-evidence & L & hash chain of entries (svc) \\
Multi-region active-active & L & conflict-free replication (infra) \\
Permission analytics & L & usage stats over checks (svc) \\
Policy as code (external engine) & M & delegate to OPA-style engine (svc) \\
Fine-grained cache keys & M & per user+resource caching (svc) \\
Schema migration framework & H & Alembic migrations (db) \\
Deadlock-safe bulk updates & M & consistent lock ordering (svc) \\
\end{longtable}

% =========================================================================
\chapter{SQL Migration Guide}
% =========================================================================

How the storage moves from Python dicts to SQLite to PostgreSQL. The
\texttt{sql\_version/} folder is the working SQLite step.

\section{Dictionary to SQLite}
\begin{itemize}[leftmargin=1.4em]
  \item Each dict becomes a table: \texttt{users}, \texttt{roles},
        \texttt{permissions}.
  \item Each ``set of ids'' becomes a many-to-many \emph{join table}:
        \texttt{user\_roles} and \texttt{role\_permissions}.
  \item The service logic barely changes --- only the storage calls do
        (a dict lookup becomes a \texttt{SELECT}; membership becomes a JOIN).
\end{itemize}

\texttt{sql\_version/} uses the standard library's \texttt{sqlite3} module
directly --- no ORM --- so every query is plain, parameterized SQL you can read
top to bottom. (If the hand-written SQL later grows unwieldy, an ORM such as
SQLAlchemy is the natural next tool; it is not needed at this size.)

\section{Schema (plain DDL)}
The whole schema is one \texttt{CREATE TABLE} script, run once at startup:
\begin{lstlisting}[language=SQL]
CREATE TABLE roles (
    id         TEXT PRIMARY KEY,
    name       TEXT NOT NULL UNIQUE,
    created_at TEXT NOT NULL,
    updated_at TEXT NOT NULL
);
CREATE TABLE role_permissions (
    role_id       TEXT    NOT NULL REFERENCES roles(id)       ON DELETE CASCADE,
    permission_id INTEGER NOT NULL REFERENCES permissions(id) ON DELETE CASCADE,
    PRIMARY KEY (role_id, permission_id)
);
\end{lstlisting}

\section{Relationships and Cascade}
Many-to-many is modelled with the two join tables. A composite primary key
(\texttt{(role\_id, permission\_id)}) blocks duplicate links. Deleting a role
removes its join rows automatically via \texttt{ON DELETE CASCADE} --- so ``delete
role'' needs no manual cleanup loop. One catch: SQLite disables foreign keys by
default, so every connection runs \texttt{PRAGMA foreign\_keys = ON}; without it
the cascade silently does nothing.

\section{JOINs}
The permission check becomes one JOIN instead of a Python loop:
\begin{lstlisting}[language=SQL]
SELECT 1
FROM user_roles ur
JOIN role_permissions rp ON rp.role_id = ur.role_id
JOIN permissions p       ON p.id = rp.permission_id
WHERE ur.user_id = ? AND p.resource = ? AND p.action = ?
LIMIT 1;
\end{lstlisting}
A returned row means ``allowed''; no row means ``denied'' (an unknown user simply
matches nothing). The \texttt{?} placeholders are bound parameters --- never
string formatting --- so the query is injection-safe.

\section{Indexes}
Index columns used in filters and joins: names, and the join-table foreign keys.
The \texttt{(resource, action)} unique constraint is itself an index. Indexes make
reads O(log n) at a small write cost.

\section{Constraints}
Push invariants into the database: \texttt{UNIQUE} on names, a composite
\texttt{UNIQUE(resource, action)}, \texttt{NOT NULL}, and foreign keys. The
database then enforces integrity even if application code has a bug. Here the app
still checks uniqueness first (for a friendly 409), but the constraint is the real
guarantee: a raced \texttt{INSERT} raises \texttt{IntegrityError}, mapped to 409.

\section{Transactions}
A request is one unit of work: several \texttt{conn.execute()} calls accumulate on
the connection, then a single \texttt{conn.commit()} makes them atomic (or
\texttt{conn.rollback()} on error). ``Replace roles'' is safe because it deletes
and re-inserts, then commits once.

\section{N+1 Query Problem}
Listing roles and loading each role's permissions with a separate query is N+1:
one query for the list plus one per role. At demo scale this is fine, and the code
keeps it that way for readability (with a comment). If the list grew, fold it into
a single JOIN and group the rows in Python:
\begin{lstlisting}[language=SQL]
SELECT r.id, p.resource, p.action
FROM roles r
LEFT JOIN role_permissions rp ON rp.role_id = r.id
LEFT JOIN permissions p       ON p.id = rp.permission_id
ORDER BY r.name;            -- one query; group by r.id in Python
\end{lstlisting}
An ORM hides this behind eager-loading options (\texttt{selectinload},
\texttt{joinedload}); with raw SQL you make the trade-off explicit.

\section{SQLite to PostgreSQL}
Swap the \texttt{sqlite3} calls for a Postgres driver (\texttt{psycopg}); the
tables and most of the SQL carry over. Add a connection pool and migrations
(Alembic) --- and this is the point where adopting an ORM starts to pay off. You
gain concurrent writers, enforced foreign keys, and replication. See
\texttt{docs/scalability-evolution.md} for the full ladder.

% =========================================================================
\chapter{System Design Evolution}
% =========================================================================

\section{Current (this repo)}
\begin{verbatim}
   Browser (frontend)
        |  HTTP + JSON
        v
   FastAPI app  ->  routes  ->  service  ->  storage
                                             (dicts OR SQLite)
\end{verbatim}
One process. State is in memory or a single SQLite file. Perfect for a demo or a
small internal tool.

\section{Production}
\begin{verbatim}
                    +----------------------+
   clients  ---->   |    Load Balancer     |
                    +----------+-----------+
                               |
              +----------------+----------------+
              v                v                v
        API node 1        API node 2       API node 3     (stateless)
              \                |                /
               \               |               /
                v              v              v
            +-------------------------------------+
            |            Redis cache              |   effective perms, O(1) reads
            +-------------------------------------+
                               |
                    +----------+-----------+
                    |    PostgreSQL primary |  <-- writes
                    +----------+-----------+
                               |  replication
                    +----------+-----------+
                    |   read replicas       |  <-- reads
                    +----------------------+
\end{verbatim}

\section{Every Layer}
\begin{itemize}[leftmargin=1.4em]
  \item \textbf{Load balancer:} spreads traffic; health-checks nodes.
  \item \textbf{API nodes:} stateless FastAPI; scale horizontally.
  \item \textbf{Cache (Redis):} serves the read-hot permission check; invalidated
        on writes.
  \item \textbf{Primary DB:} source of truth for writes; transactions and
        constraints.
  \item \textbf{Replicas:} absorb read traffic; accept small lag.
  \item \textbf{(Optional) event bus + workers:} async cache invalidation, audit
        fan-out, notifications.
\end{itemize}

% =========================================================================
\chapter{Interview Memorization Guide}
% =========================================================================

\section{Remember First (in order)}
1. Layering: routes \texttt{->} service \texttt{->} storage. 2. The check is a set
membership: \texttt{Permission(resource, action) in effective\_permissions}.
3. References not copies (users hold role ids). 4. Fail closed / default deny.
5. Validate before mutate; all-or-nothing.

\section{FastAPI Syntax}
\begin{lstlisting}[language=Python]
app = FastAPI(); router = APIRouter()
@router.post("/users", response_model=UserResponse, status_code=201)
def create_user(body: CreateUserRequest, svc = Depends(get_service)): ...
@router.get("/users/{user_id}", response_model=UserResponse)
def get_user(user_id: str, name: str | None = Query(default=None)): ...
app.include_router(router)
@app.exception_handler(AppError)
def handle(request, exc): return JSONResponse(exc.status_code, {...})
\end{lstlisting}

\section{Pydantic Syntax}
\begin{lstlisting}[language=Python]
class CreateUserRequest(BaseModel):
    name: str = Field(min_length=1)
    role_ids: list[str] = Field(default_factory=list)
# validate: CreateUserRequest(**payload)   serialize: model.model_dump()
\end{lstlisting}

\section{sqlite3 Syntax}
\begin{lstlisting}[language=Python]
conn = sqlite3.connect(path); conn.row_factory = sqlite3.Row
conn.execute("PRAGMA foreign_keys = ON")            # enable cascades
conn.execute("SELECT * FROM users WHERE id = ?", (uid,)).fetchone()   # one row
conn.execute("SELECT * FROM users").fetchall()      # many rows
conn.execute("INSERT INTO users(...) VALUES(?,...)", vals); conn.commit()
conn.execute("DELETE FROM roles WHERE id = ?", (rid,)); conn.commit()  # + cascade
\end{lstlisting}

\section{CRUD Patterns}
\begin{lstlisting}[language=Python]
# Create: validate -> construct -> add -> commit -> return
# Read:   get by id (404 if None) OR list with filters
# Update: fetch -> validate -> mutate fields -> commit
# Delete: fetch (404) -> delete -> commit (handle cascade)
\end{lstlisting}

\section{Dependency Injection}
\begin{lstlisting}[language=Python]
def get_service(): return service          # or: def get_db(): yield connect()
def endpoint(svc = Depends(get_service)): ...
app.dependency_overrides[get_service] = lambda: test_service   # in tests
\end{lstlisting}

\section{REST Patterns}
GET read, POST create, PUT replace, PATCH partial, DELETE remove. 200 ok, 201
created, 204 no content, 400/422 bad input, 401/403 auth, 404 missing, 409
conflict.

\section{Error Handling}
Raise typed exceptions in the service (\texttt{NotFoundError},
\texttt{ConflictError}, \texttt{ValidationError}); one handler per type in
\texttt{main.py} renders a consistent JSON body. Never scatter try/except across
routes.

% =========================================================================
\chapter{Mock Interview: 75 Questions}
% =========================================================================

Each question has an ideal answer, hints, common mistakes, and a follow-up.

\section{Coding \& Implementation}
\mq{Implement delete user. Walk me through it.}
{Add \texttt{delete\_user} in the service (validate exists, pop from store,
invalidate cache), then a DELETE route returning 204.}
{Where does the logic go? What status code?}
{Doing it in the route; returning 200 with a body; not handling missing id.}
{Now make role deletion cascade.}

\mq{How do you delete a role without leaving dangling references?}
{Delete the role, then discard its id from every user's role set (in memory) or
let the join table cascade (SQL).}
{Users hold role ids --- what happens to them?}
{Forgetting the cascade; leaving stale ids that break later reads.}
{What is the complexity of the cascade?}

\mq{Add pagination to list users.}
{Accept \texttt{limit} and \texttt{offset} query params; slice in memory or
\texttt{LIMIT/OFFSET} in SQL; return metadata.}
{Do you want offset or cursor pagination?}
{Loading all rows then slicing in Python at scale; no upper bound on limit.}
{Why is cursor pagination better for large tables?}

\mq{Implement bulk-assign roles to a user.}
{Validate every role id first (all-or-nothing), then apply, then invalidate
cache; one transaction in SQL.}
{What if one id is invalid halfway through?}
{Half-applying and leaving inconsistent state.}
{How do you make it idempotent?}

\mq{Add rename role.}
{PATCH route; service trims, checks uniqueness excluding self, updates name and
\texttt{updated\_at}.}
{PATCH or PUT? What about the current name?}
{Uniqueness check that does not exclude the record being renamed.}
{How would you keep a rename history?}

\mq{Implement effective permissions for a user.}
{Union the permission sets of all the user's roles; cache it.}
{Multiple roles --- how do they combine?}
{Returning the first role's perms only; not de-duplicating.}
{How do you keep the cache correct?}

\mq{Add search by name.}
{Query param, case-insensitive substring filter over the collection.}
{New path or query param?}
{Creating a new endpoint instead of a filter; case-sensitive match.}
{Add pagination and sorting on top.}

\mq{Add wildcard permissions like documents:*.}
{On check, test the exact permission and any wildcard forms in the effective
set.}
{How do you match without slowing the check?}
{Letting wildcards silently over-grant; slow scans.}
{How do explicit denies interact with wildcards?}

\mq{Implement soft delete.}
{Add \texttt{deleted\_at}; deletes set it; all reads filter it out.}
{Hard or soft --- what does the product need?}
{Forgetting to filter soft-deleted rows in some query.}
{How do you support restore and GDPR erase?}

\mq{Add created\_at and updated\_at.}
{Set both on create; set \texttt{updated\_at} only on real modifications.}
{When should updated\_at change?}
{Bumping updated\_at on reads; timezone confusion.}
{Store UTC or local? Why UTC?}

\mq{Return a consistent error format.}
{Typed exceptions plus one handler that emits
\texttt{\{error:\{code,message\}\}}.}
{Where do you centralize this?}
{try/except in every route; inconsistent shapes.}
{Add a correlation id to each error.}

\mq{Implement replace-permissions for a role.}
{PUT; normalize and set the whole permission set; invalidate cache.}
{Replace or merge?}
{Merging when replace was asked; not normalizing case.}
{How to audit what changed?}

\mq{Add a health endpoint.}
{Return \texttt{status: ok} without touching storage so it is fast and always
answers.}
{Should it query the DB?}
{Making health depend on the DB so a slow DB fails health.}
{Liveness vs readiness?}

\mq{Add counts for users and roles.}
{\texttt{len} in memory or \texttt{SELECT count(*)} in SQL; simple response
schema.}
{Static path ordering vs \texttt{/users/\{id\}}?}
{\texttt{/users/count} captured by \texttt{/users/\{user\_id\}} due to order.}
{How to count users per role efficiently?}

\mq{Implement an audit log.}
{Append an entry (timestamp, action, detail) on each write; expose a read
endpoint.}
{Where do you record it?}
{Mutating past entries; unbounded growth.}
{How do you make it tamper-evident?}

\section{Architecture}
\mq{Why keep logic in the service, not the route?}
{Testability and reuse; routes stay thin adapters; storage can change without
touching rules.}
{What would break if logic lived in routes?}
{Business logic entangled with HTTP; hard to unit test.}
{Show a test that hits the service directly.}

\mq{Why do users store role ids, not role objects?}
{Single source of truth; editing a role affects all holders instantly.}
{What happens on a role edit?}
{Embedding copies that drift out of sync.}
{What is the downside of references?}

\mq{How would you swap storage for a database?}
{Implement the same method surface over a DB; the service is unchanged.}
{What is the seam?}
{Leaking SQL into the service.}
{Show where the in-memory and SQL services differ.}

\mq{Where does validation belong?}
{At the edge (Pydantic) for shape, and in the service for rules
(uniqueness, references).}
{Two layers --- which does what?}
{Trusting client input; duplicating checks inconsistently.}
{What can Pydantic not check?}

\mq{How do you keep responses deterministic?}
{Sort collections (role ids, permissions) before returning.}
{Why might output order vary?}
{Relying on set iteration order in tests.}
{Why does this help testing?}

\mq{Explain the layering top to bottom.}
{Frontend -> route (validate/format) -> service (rules) -> storage (data) and
back.}
{What is each layer responsible for?}
{Blurring responsibilities across layers.}
{Where would auth middleware sit?}

\mq{How do you add authentication later without a rewrite?}
{A dependency that verifies a token and yields the user; protect routes with it;
service unchanged.}
{Which layer changes?}
{Trusting a user id from the body.}
{Access vs refresh tokens?}

\mq{How would you support multi-tenancy?}
{Add a tenant id, scope every query and uniqueness check by it.}
{What becomes the effective key?}
{Global uniqueness leaking across tenants.}
{Where does the tenant id come from at request time?}

\section{FastAPI \& Python}
\mq{What does Depends do?}
{Injects a dependency's return value; enables shared setup and test overrides.}
{How do you replace it in tests?}
{Using a bare global that cannot be swapped.}
{Can dependencies be nested?}

\mq{Why declare static routes before parametrized ones?}
{First-match wins; else \texttt{/users/count} is captured by
\texttt{/users/\{user\_id\}}.}
{What matches \texttt{/users/count} first?}
{Wrong order causing a 404 on count.}
{How does Starlette resolve routes?}

\mq{How does response\_model help?}
{Serializes and filters output to a declared shape; documents the endpoint.}
{What if you return extra internal fields?}
{Leaking internals; shape drift.}
{Can it hide secrets?}

\mq{What is the difference between PUT and PATCH here?}
{PUT replaces the whole collection (roles); PATCH updates part (name).}
{Which did you use for replace-roles vs rename?}
{Using PATCH to replace or vice versa.}
{Which is idempotent?}

\mq{Why a frozen dataclass for Permission?}
{Immutable means hashable, so it lives in a set with O(1) membership and
de-dup.}
{What error do you get otherwise?}
{Unhashable type when adding to a set.}
{What does frozen generate?}

\mq{Why default\_factory for list/set fields?}
{A mutable default is shared across instances; the factory makes a fresh one.}
{What is the classic bug?}
{Shared mutable default causing cross-instance state.}
{Show the trap.}

\mq{dict.get vs indexing?}
{\texttt{get} returns None on miss; \texttt{[]} raises KeyError. Use get to fail
closed.}
{When do you want the exception?}
{Using \texttt{[]} where absence is expected.}
{Where is this used in the check?}

\mq{How does FastAPI produce a 422?}
{Pydantic validation fails before your code runs; a handler can reshape it.}
{What triggers it here?}
{Manually validating inside handlers.}
{What does the error body contain?}

\section{REST \& API Design}
\mq{Which status code for create, update, delete?}
{201, 200, 204 respectively; 404/409/422 for errors.}
{Why 201 not 200 on create?}
{200 on create; 200 on delete with a body.}
{What should a 201 return?}

\mq{Is your API idempotent where it should be?}
{GET/PUT/DELETE and set-based POSTs are; plain create is not.}
{Which operations repeat safely?}
{Assuming all POSTs are idempotent.}
{How to make create idempotent?}

\mq{Body on a DELETE --- is that allowed?}
{Yes; we send the permission to remove in the body since it has no id.}
{How else could you identify it?}
{Assuming DELETE cannot carry a body.}
{Encode it in the path instead?}

\mq{How do you version an API?}
{Path prefix (/v1) or headers; keep old versions until clients migrate.}
{Where would you add the prefix?}
{Breaking clients with silent changes.}
{URL vs header versioning trade-offs?}

\mq{How would you paginate and still return a total?}
{Return items plus metadata (total, next cursor).}
{Offset or cursor?}
{Counting on every page for huge tables.}
{Why can total be expensive?}

\mq{What makes an endpoint RESTful?}
{Resource URLs, standard verbs, statelessness, correct status codes.}
{Is this API fully REST?}
{GET that mutates state.}
{Where is it pragmatic rather than pure REST?}

\section{SQL}
\mq{Model the many-to-many relationships.}
{Join tables user\_roles and role\_permissions with composite primary keys and
foreign keys.}
{How many join tables?}
{Storing role ids in a column instead of a join table.}
{Why composite primary keys?}

\mq{Write the permission check as SQL.}
{JOIN permissions to role\_permissions to user\_roles, filter by user id and the
resource/action, limit 1.}
{Which tables join?}
{Loading everything into Python and filtering there.}
{How do indexes help this query?}

\mq{What is the N+1 problem and how do you fix it?}
{Loading each row's children with a separate query fires one query per row in a
loop; fix by folding it into a single JOIN and grouping in Python (an ORM would
instead use eager loading like selectinload).}
{What triggers the extra queries?}
{A per-row query inside a list loop.}
{How would an ORM express the same fix?}

\mq{How does delete-role cascade in SQL?}
{A DELETE on the role removes its join-table rows via ON DELETE CASCADE, so no
dangling links --- provided PRAGMA foreign\_keys = ON is set on the connection.}
{Do you loop over users?}
{Forgetting to enable foreign keys, so the cascade silently does nothing.}
{How would you verify the join rows are actually gone?}

\mq{Where do you add indexes and why?}
{On names and join-table foreign keys and (resource, action); they speed reads at
a small write cost.}
{Index everything?}
{Over-indexing; slow writes, wasted space.}
{Which query does each index serve?}

\mq{What does a transaction guarantee?}
{Atomicity: a group of writes all succeed or all roll back.}
{Where do you need it here?}
{Multi-step writes committed separately.}
{Why keep transactions short?}

\mq{Lazy vs eager loading (and the raw-SQL equivalent)?}
{In an ORM, lazy loads related rows on access (convenient, N+1 risk) and eager
loads them up front (use for lists). With raw sqlite3 there is no lazy magic ---
you either issue a query per row or write one JOIN, so the trade-off is explicit.}
{Which for a list endpoint?}
{A per-row query in a loop.}
{How does raw SQL make the choice visible?}

\mq{How would you migrate SQLite to Postgres?}
{Swap sqlite3 for a Postgres driver (psycopg); most SQL carries over. Add a
connection pool and migrations (Alembic), and consider an ORM as the SQL grows.}
{What do you gain?}
{Assuming a full rewrite is needed.}
{What new concerns appear (pooling, concurrency)?}

\section{RBAC \& Security}
\mq{Explain default deny.}
{No matching permission means denied; access is never implicit.}
{What does an unknown user get?}
{Failing open on errors.}
{Where do you fail closed?}

\mq{How do wildcard and explicit-deny interact?}
{Deny should override allow; evaluate denies after allows.}
{What wins in a conflict?}
{Allow silently overriding a deny.}
{How do you keep this fast?}

\mq{Why is trusting user\_id from the body insecure?}
{Anyone could impersonate; identity must come from a verified token.}
{Where should identity come from?}
{Shipping this as production auth.}
{How does JWT fix it?}

\mq{How would you enforce least privilege?}
{Design narrow roles; grant only what a job needs; review regularly.}
{One super-role for everyone?}
{Over-broad roles.}
{How do you find over-privileged users?}

\mq{Why audit logs matter?}
{They record who changed or checked access --- essential for compliance and
incident response.}
{What fields do you log?}
{No trail after an incident.}
{How to make them tamper-evident?}

\mq{What is separation of duties?}
{Constraints preventing one person from holding conflicting roles.}
{Give an example.}
{Ignoring conflicting-role rules.}
{How would you model it?}

\mq{Risks of allow-all CORS?}
{Any site can call the API from a browser; fine locally, unsafe with credentials
in production.}
{Why is it fine here?}
{Leaving it open in production.}
{How do you scope origins?}

\section{Testing}
\mq{What do you test for a new feature?}
{Happy path, not-found, validation error, duplicate, and idempotency.}
{Which cases are non-negotiable?}
{Only testing the happy path.}
{Unit vs integration here?}

\mq{How do you isolate tests?}
{Fresh service/DB per test via fixtures; override dependencies.}
{Why fresh state?}
{Tests sharing state and order-dependence.}
{How do you inject a clean store?}

\mq{How do you test an error is raised?}
{\texttt{with pytest.raises(NotFoundError): ...} at the service; assert status
codes at the API.}
{Service or API level?}
{try/except that never fails if no error.}
{How to assert the message/code?}

\mq{How do you isolate the SQL version's tests?}
{Point get\_db at a fresh temp-file SQLite DB per test and override the
dependency; do not enter TestClient as a context manager, so the app's startup
lifespan never touches the real rbac.db.}
{Why a temp file and not :memory:?}
{With a new connection per request, :memory: gives each request its own empty
database; a temp file is shared across them.}
{How do you reset between tests?}

\mq{How do you test cache invalidation?}
{Compute perms (populate cache), change a permission, assert the check reflects
it.}
{What proves the cache was invalidated?}
{Only checking before the change.}
{Which test covers this?}

\mq{What is a regression test?}
{A test added with a bug fix so the bug cannot return.}
{When do you add one?}
{Fixing without a test.}
{Give an example here.}

\mq{Why both service and API tests?}
{Service tests are fast unit tests; API tests verify routing, validation, and
serialization.}
{What does each catch?}
{Calling everything a unit test.}
{Which catches a wrong status code?}

\section{Scalability}
\mq{First step to scale past one process?}
{Externalize state to a shared DB so stateless nodes agree.}
{What breaks with in-memory at scale?}
{Scaling nodes while keeping in-memory state.}
{Then what for the read-hot check?}

\mq{How does caching change the check?}
{Cache the effective permission set; a check is an O(1) lookup.}
{What is the hard part?}
{Never invalidating.}
{When do you invalidate?}

\mq{Cache invalidation strategy for a role edit?}
{Invalidate every user holding that role (fan-out) or clear broadly; TTL as a
safety net.}
{Who is affected by a role change?}
{Invalidating only one user.}
{How do you bound the fan-out?}

\mq{How do read replicas help and hurt?}
{Reads scale out; the cost is replication lag (eventual consistency).}
{Where do writes go?}
{Reading your own write from a lagging replica.}
{Which reads must hit the primary?}

\mq{Role hierarchy at scale?}
{Precompute the transitive closure and cache it; recompute on hierarchy edits.}
{Why not compute per request?}
{Traversing the graph on every check.}
{How do you detect cycles?}

\mq{How do you answer "who can do X" at scale?}
{Maintain a reverse index permission -> roles -> users, updated on writes.}
{Why not scan?}
{Full scans per query.}
{Where is that index stored?}

\mq{What is optimistic locking?}
{A version column; update only if the version matches, else 409 and retry.}
{Why not lock rows?}
{Lost updates from concurrent edits.}
{Optimistic vs pessimistic trade-offs?}

\mq{Where is the likely bottleneck and how do you confirm?}
{The permission check on the read path; confirm with metrics/profiling, then
cache.}
{Read or write path?}
{Guessing without data.}
{What metric signals trouble?}

\section{Rapid Fire}
\mq{Why return 204 and not 200 on delete?}
{204 means success with no body, which fits a delete that has nothing to
return.}
{What does the client need back?}
{Returning 200 with a body or an empty 200.}
{When would you return the deleted object instead?}

\mq{How do you make list output deterministic for tests?}
{Sort by a stable key (id or name) before returning.}
{Why can order vary?}
{Relying on dict/set iteration order.}
{Where do the converters sort?}

\mq{Add a "who am I" endpoint returning a user and their effective permissions.}
{Combine get\_user and get\_effective\_permissions into one response.}
{Reuse existing service methods?}
{Recomputing perms without the cache.}
{How would auth supply the user id here?}

\mq{How would you prevent deleting a role that is still in use?}
{Count holders; if greater than zero, return 409 instead of cascading.}
{Cascade or block --- product decision?}
{Silently cascading when the product wanted a block.}
{How do you report the blocking holders?}

\mq{How do you keep the SQL and in-memory versions behavior-compatible?}
{Share the same rules and run the same API-level tests against both.}
{What differs and what stays the same?}
{Letting the two drift in validation or status codes.}
{Could you extract a common test suite?}

\mq{What breaks if you store role objects inside users in SQL?}
{Duplication and update anomalies; the join table is the normalized answer.}
{Why normalize?}
{Denormalizing without a reason.}
{When is denormalization justified?}

\mq{How do you validate the permission vocabulary (allowed resources/actions)?}
{Check against an allow-list in the service; reject unknown values with 422.}
{Where does this belong?}
{Accepting any string, leading to typos that never match.}
{How do you evolve the vocabulary safely?}

\mq{How would you add pagination metadata without an expensive count?}
{Return items plus a next cursor; skip total, or cache/estimate it.}
{Why is count(*) costly on large tables?}
{Counting the whole table on every page.}
{Cursor vs offset for this?}

% =========================================================================
\chapter{Whiteboard Exercises (approx. 40)}
% =========================================================================

For each: the problem, the expected approach, files to modify, pseudocode, and
complexity. Try each before reading the approach.

\section{Core CRUD and Membership}

\textbf{WB1. Delete user.} Files: svc, route.
\begin{lstlisting}
def delete_user(user_id): get(404); pop; invalidate_cache(user_id)
route: DELETE /users/{id} -> 204
\end{lstlisting}
Complexity: O(1).

\textbf{WB2. Delete role with cascade.} Files: svc, route.
\begin{lstlisting}
def delete_role(role_id):
    get(404); del roles[role_id]
    for u in users: u.role_ids.discard(role_id)
    invalidate_cache()
\end{lstlisting}
Complexity: O(U).

\textbf{WB3. Get user by id.} Files: route. \texttt{get\_user} then 404 if None.
Complexity: O(1).

\textbf{WB4. Rename role.} Files: sch, svc, route.
\begin{lstlisting}
def rename_role(id, name):
    role = get(id); name = clean(name)
    assert_unique(name, exclude_id=id)
    role.name = name; role.updated_at = now()
\end{lstlisting}
Complexity: O(N) memory / O(1) SQL with unique index.

\textbf{WB5. Replace user roles.} Files: sch, svc, route.
\begin{lstlisting}
def replace_user_roles(id, role_ids):
    u = get(id)
    for r in role_ids: get_role(r)   # validate all first
    u.role_ids = set(role_ids); invalidate_cache(id)
\end{lstlisting}
Complexity: O(K).

\textbf{WB6. Effective permissions.} Files: svc, route.
\begin{lstlisting}
def effective(id):
    if id in cache: return cache[id]
    perms = union(role.permissions for role in user's roles)
    cache[id] = perms; return perms
\end{lstlisting}
Complexity: O(R*P), O(1) cached.

\textbf{WB7. Users having a role.} Files: svc, route.
\begin{lstlisting}
def users_with_role(rid): get_role(rid); return [u for u in users if rid in u.role_ids]
\end{lstlisting}
Complexity: O(U) / O(matches) with reverse index.

\textbf{WB8. Count users per role.} Files: svc, route.
\begin{lstlisting}
counts = {r.id: 0 for r in roles}
for u in users:
    for rid in u.role_ids: counts[rid] += 1
\end{lstlisting}
Complexity: O(U*R\_avg).

\section{Search, Pagination, Sorting}

\textbf{WB9. Search users by name.} Files: svc, route.
\begin{lstlisting}
def search(name): q = name.strip().lower(); return [u for u in users if q in u.name.lower()]
\end{lstlisting}
Complexity: O(N).

\textbf{WB10. Offset pagination.} Files: svc, route.
\begin{lstlisting}
route: GET /users?limit=&offset=
return items[offset:offset+limit], total=len(items)
\end{lstlisting}
Complexity: O(N) memory / O(limit) SQL with index.

\textbf{WB11. Cursor pagination.} Files: svc, route.
\begin{lstlisting}
route: GET /users?after=<id>&limit=
SQL: WHERE id > :after ORDER BY id LIMIT :limit
\end{lstlisting}
Complexity: O(limit) with an index on id.

\textbf{WB12. Sort by created\_at.} Files: svc, route.
\begin{lstlisting}
key = query param; reverse = (order == "desc")
return sorted(items, key=attrgetter(key), reverse=reverse)
\end{lstlisting}
Complexity: O(N log N).

\textbf{WB13. Multi-filter (name + has-role).} Files: svc, route.
\begin{lstlisting}
items = users
if name: filter by substring
if role_id: filter role_id in u.role_ids
\end{lstlisting}
Complexity: O(N).

\section{Bulk and Import/Export}

\textbf{WB14. Bulk assign roles.} Files: svc, route.
\begin{lstlisting}
def bulk_assign(user_id, role_ids):
    u = get(user_id)
    for r in role_ids: get_role(r)   # validate all
    u.role_ids |= set(role_ids); invalidate_cache(user_id)
\end{lstlisting}
Complexity: O(K).

\textbf{WB15. Bulk create users (transaction).} Files: svc, route.
\begin{lstlisting}
validate all names unique + non-blank
then create all; SQL: single transaction, commit once
\end{lstlisting}
Complexity: O(K).

\textbf{WB16. Export all to JSON.} Files: svc, route.
\begin{lstlisting}
return {"users": [to_dict(u)...], "roles": [to_dict(r)...]}
\end{lstlisting}
Complexity: O(N).

\textbf{WB17. Import from JSON.} Files: svc, route.
\begin{lstlisting}
validate every record; create roles first, then users referencing them
report per-record success/failure
\end{lstlisting}
Complexity: O(N).

\section{Permissions Model Extensions}

\textbf{WB18. Wildcard permission check.} Files: svc.
\begin{lstlisting}
candidates = {Permission(res, act), Permission(res, "*"), Permission("*", act)}
return any(c in effective(user) for c in candidates)
\end{lstlisting}
Complexity: O(1) extra lookups.

\textbf{WB19. Explicit deny overrides allow.} Files: models, svc.
\begin{lstlisting}
if required in denies(user): return False
return required in allows(user)
\end{lstlisting}
Complexity: O(1) with sets.

\textbf{WB20. Role hierarchy (inheritance).} Files: models, svc.
\begin{lstlisting}
def effective_perms(role_id, seen=set()):
    if role_id in seen: return set()   # cycle guard
    seen.add(role_id); perms = role.permissions
    for parent in role.parents: perms |= effective_perms(parent, seen)
    return perms
\end{lstlisting}
Complexity: O(V+E); cache the result.

\textbf{WB21. Permission groups.} Files: models, svc.
\begin{lstlisting}
group "content" = {documents:read, documents:write}
on assign: expand group into the role's permission set
\end{lstlisting}
Complexity: O(size of group).

\textbf{WB22. Time-boxed role assignment.} Files: models, svc.
\begin{lstlisting}
store (role_id, until); effective() skips assignments where until < now
\end{lstlisting}
Complexity: O(R).

\textbf{WB23. Batch permission check.} Files: svc, route.
\begin{lstlisting}
perms = effective(user)         # once
return {f"{r}:{a}": Permission(r,a) in perms for (r,a) in requests}
\end{lstlisting}
Complexity: O(K).

\textbf{WB24. Field-level permissions.} Files: svc.
\begin{lstlisting}
permission = (resource, action, field); check includes field
\end{lstlisting}
Complexity: O(1) with sets.

\section{Caching, Concurrency, Reliability}

\textbf{WB25. Effective-permission cache with invalidation.} Files: store, svc.
\begin{lstlisting}
read: return cache[user] or compute+store
write (role change): invalidate affected users
\end{lstlisting}
Complexity: O(1) hit; O(R*P) miss.

\textbf{WB26. Invalidate all holders on role edit.} Files: svc.
\begin{lstlisting}
def on_role_perm_change(role_id):
    for u in users_with_role(role_id): invalidate_cache(u.id)
\end{lstlisting}
Complexity: O(holders).

\textbf{WB27. Optimistic locking.} Files: models, svc.
\begin{lstlisting}
UPDATE role SET ..., version=version+1 WHERE id=? AND version=?
if rowcount == 0: raise ConflictError   # someone else changed it
\end{lstlisting}
Complexity: O(1).

\textbf{WB28. Idempotency key on create.} Files: store, route.
\begin{lstlisting}
if key in seen: return seen[key]
result = create(...); seen[key] = result; return result
\end{lstlisting}
Complexity: O(1).

\textbf{WB29. Rate limiter (token bucket).} Files: middleware.
\begin{lstlisting}
tokens[client] refill by rate*elapsed, cap at burst
if tokens < 1: 429 else tokens -= 1
\end{lstlisting}
Complexity: O(1) per request.

\textbf{WB30. Soft delete + filter.} Files: models, svc.
\begin{lstlisting}
delete: set deleted_at = now
all reads: WHERE deleted_at IS NULL
\end{lstlisting}
Complexity: O(1) write; reads need a partial index.

\textbf{WB31. Audit log with actor.} Files: models, svc.
\begin{lstlisting}
on each write: append {ts, actor, action, target, detail}  (append-only)
\end{lstlisting}
Complexity: O(1) per write.

\textbf{WB32. Cache stampede guard.} Files: svc.
\begin{lstlisting}
on miss: acquire per-key lock; recompute once; others wait/read
\end{lstlisting}
Complexity: O(1) amortized.

\section{SQL-specific}

\textbf{WB33. Permission check as one JOIN.} Files: svc (SQL).
\begin{lstlisting}
select 1 from permissions p
join role_permissions rp on p.id = rp.permission_id
join user_roles ur on rp.role_id = ur.role_id
where ur.user_id=? and p.resource=? and p.action=? limit 1
\end{lstlisting}
Complexity: O(log n) with indexes.

\textbf{WB34. Avoid N+1 in list users.} Files: svc (SQL).
\begin{lstlisting}
select u.id, ur.role_id from users u
left join user_roles ur on ur.user_id = u.id
order by u.name        -- one query; group role_ids by u.id in Python
\end{lstlisting}
Complexity: 1 query instead of N+1.

\textbf{WB35. Users-with-role via JOIN.} Files: svc (SQL).
\begin{lstlisting}
select u.* from users u
join user_roles ur on ur.user_id = u.id
where ur.role_id = ?
\end{lstlisting}
Complexity: O(matches) with an index on role\_id.

\textbf{WB36. Count users per role (GROUP BY).} Files: svc (SQL).
\begin{lstlisting}
select role_id, count(*) from user_roles group by role_id
\end{lstlisting}
Complexity: O(rows) with index.

\textbf{WB37. Unique constraint on permission pair.} Files: models (SQL).
\begin{lstlisting}
UniqueConstraint("resource", "action")  -> get-or-create avoids duplicates
\end{lstlisting}
Complexity: O(1) lookup on the index.

\textbf{WB38. Zero-downtime add-column migration.} Files: migrations.
\begin{lstlisting}
add nullable column -> backfill in batches -> make not-null -> deploy code
\end{lstlisting}
Complexity: backfill O(N) in batches.

\section{System Design Sketches}

\textbf{WB39. Add Redis cache in front of checks.} Files: svc, infra.
\begin{lstlisting}
check: GET redis user:perms; on miss compute from DB and SETEX
invalidate on role/permission writes (publish event)
\end{lstlisting}
Complexity: O(1) hit; network on miss.

\textbf{WB40. Central authorization service.} Files: new service.
\begin{lstlisting}
other services call POST /check-permission (or gRPC)
authz service owns roles/perms; caches; scales independently
\end{lstlisting}
Complexity: one network hop per check (cache to keep it cheap).

\vspace{1em}
\begin{center}
\textit{End of feature-extensions guide. See \texttt{report.tex} for the core
project report and \texttt{docs/scalability-evolution.md} for the full scaling
narrative.}
\end{center}

\end{document}
